\section{From plane curves to normal polynomials}
\label{sec:reductions}

We reduce a violation at very general points to an equal-multiplicity
system existing at every configuration.  Moving its points toward a
smooth plane curve then gives a polynomial on the normal bundle.
This polynomial will exclude the square cases and supply the input
for the cubic construction in Section~\ref{sec:elliptic}.

The order of a section is the order of its scalar expression in a local
frame.  A frame change multiplies this expression by a unit, preserving
the order.  In local coordinates, order at least $m$ means that every
Taylor coefficient of total degree less than $m$ vanishes.  Thus
algebraic and holomorphic orders agree, and holomorphic coordinate
changes preserve them.

\subsection{A universal equal-multiplicity system}

For positive integers $d,m$, write $\mathsf U_r(d,m)$ for the property
\begin{equation}
\label{eq:universal}
\begin{gathered}
\text{For every }(p_1,\ldots,p_r)\in U_r\text{ there is a nonzero form}\\
S\in H^0(\Pp^2,\mathcal O_{\Pp^2}(d))
\quad\text{such that}\quad
\mult_{p_i}(S)\ge m\quad(1\le i\le r).
\end{gathered}
\end{equation}
Here and below the multiplicity of a form is the multiplicity of its
zero divisor.

\begin{lemma}[Reduction to universal systems]
\label{lem:universal}
For each $r\ge10$ there is a countable union $E_r$ of proper
Zariski-closed subsets of $U_r$, all defined over $\mathbb Q$, with
$U_r\setminus E_r\ne\varnothing$, having the following property.
If a tuple outside $E_r$ supports a
degree-$d_0$ form, $d_0\ge1$, with multiplicities at least
$m_1,\ldots,m_r\in\Z_{\ge0}$ and
\[
 d_0\sqrt r\le\sum_{i=1}^r m_i,
\]
then $\mathsf U_r(d,m)$ holds for some positive integers $d,m$ with
$d/m\le\sqrt r$.  If $\mathsf U_r(d,m)$ holds, then
$\mathsf U_r(Nd,Nm)$ holds for every positive integer $N$.
Consequently, to prove Theorem~\ref{thm:main} for a fixed $r$, it suffices
to rule out every $\mathsf U_r(d,m)$ with $d/m\le\sqrt r$.
\end{lemma}

\begin{proof}
For a positive integer $e$ and
$\mathbf n=(n_1,\ldots,n_r)\in\Z_{\ge0}^r$, let
$A_{e,\mathbf n}\subseteq U_r$ consist of the tuples admitting a nonzero
degree-$e$ form with multiplicity at least $n_i$ at the $i$th point.
In standard affine charts, the required Taylor coefficients give a
matrix with $\binom{e+2}{2}$ columns, one for each form coefficient.
Its entries are polynomials with integer coefficients in the point
coordinates.  A nonzero form exists exactly when the matrix has less
than full column rank.  Thus $A_{e,\mathbf n}$ is cut out by its
full-column-size minors; if there are fewer rows than columns, the
condition is automatic.  Invariance of order makes these rank
conditions agree on overlaps.  Consequently $A_{e,\mathbf n}$ is
Zariski closed and defined over $\mathbb Q$.

Let $E_r$ be the union of all proper $A_{e,\mathbf n}$.  This is a
countable union.  On the dense chart where all points are affine,
each proper $A_{e,\mathbf n}$ has a defining minor that is a nonzero
polynomial over $\mathbb Q$.  Choose
$2r$ complex numbers algebraically independent over $\mathbb Q$ as
affine coordinates of the marked points.  Such numbers can be chosen
successively, since the algebraic closure of a countable field is
countable.  None of these minors vanishes at the resulting tuple, which
also avoids the diagonals.  Hence it lies in $U_r\setminus E_r$.

Now suppose that a violation as in the statement occurs outside $E_r$.
Then $A_{d_0,\mathbf m}=U_r$.  We use the cyclic-product reduction of
Nagata~\cite{Nagata1965}.
Permuting the ordered points gives the
same assertion for each cyclic shift of $\mathbf m$.  At any fixed
tuple, choose a nonzero form for each of the $r$ cyclic shifts and
take their product.  This product is nonzero, and orders add because
the product of the initial Taylor terms is nonzero.  Each original
multiplicity occurs once at each point.  Hence the product has degree
and common lower multiplicity
\[
 d=rd_0,
 \qquad
 m=\sum_{i=1}^r m_i.
\]
The assumed violation makes $m>0$, and
\[
 \frac d m
 =\frac{r d_0}{\sum_i m_i}
 \le\sqrt r.
\]
The construction applies at every tuple, proving $\mathsf U_r(d,m)$;
taking $N$th powers proves the common-multiple assertion.  Thus
excluding these universal systems excludes every violation outside
the same $E_r$, for all degrees and multiplicity vectors.
\end{proof}

\subsection{Specialization to a normal bundle}

Let $B\subseteq\Pp^2$ be a smooth irreducible plane curve of degree
$k\ge1$, and fix a nonzero homogeneous equation $G_B$ for it.  Put
\[
 L=\mathcal O_B(1),\qquad M=L^k.
\]
The equation identifies the normal bundle of $B$ in $\Pp^2$ with $M$,
as in the formula for an effective Cartier divisor
\cite[Tag~0B3P]{Stacks}.  In a local frame $e$ for
$\mathcal O_{\Pp^2}(1)$, the local equation
$t=G_B/e^k$ identifies a normal vector with its derivative under $t$.
For $w\in M_a$, the expression $f_j(a)w^j$ is an element of $L_a^d$
when $f_j$ is a section of $L^{d-kj}$.  Thus the polynomial in the next
proposition is intrinsically a section of the pullback of $L^d$ to
$\Tot(M)$.

\begin{proposition}[A necessary normal polynomial]
\label{prop:normal-polynomial}
Suppose that $\mathsf U_r(d,m)$ holds.  For every choice of distinct
points $a_1,\ldots,a_r\in B$ and arbitrary vectors $c_i\in M_{a_i}$,
there is a nonzero polynomial
\begin{equation}
\label{eq:normal-polynomial}
 F(w)=\sum_{j=0}^{\lfloor d/k\rfloor}f_jw^j,
 \qquad f_j\in H^0(B,L^{d-kj}),
\end{equation}
having multiplicity at least $m$ at each $(a_i,c_i)\in\Tot(M)$.
Its coefficients may be taken to be algebraic sections.
\end{proposition}

\begin{proof}
We move the marked points in the plane, rescale the normal coordinate,
and retain the first nonzero term of a family of plane forms.

Choose a plane affine chart and its standard frame near each $a_i$.
The differential of the local equation $t_i$ is nonzero by smoothness,
so choose an affine tangent vector $V_i$ with $dt_i(V_i)=c_i$ in
that frame.  The affine-linear motions
\[
 p_i(s)=a_i+sV_i
\]
are rational in $s$, stay distinct for sufficiently small $s$, and satisfy
\[
 \lim_{s\to0}\frac{t_i(p_i(s))}{s}=c_i
\]
in the chosen frames.  The vectors $c_i$ are arbitrary complex
normal vectors; no compatibility among them is required.

Multiplicity at least $m$ at all $p_i(s)$ gives a homogeneous linear
system over $\C(s)$ in the degree-$d$ form coefficients.  It has a
nonzero solution: otherwise some full-column-size minor would be a
nonzero rational function, giving full column rank at general complex
values of $s$ and contradicting $\mathsf U_r(d,m)$.

Clear denominators in a nonzero rational solution, and divide its
polynomial coefficients by their maximal common power of $s$.  This
gives a finite polynomial family of homogeneous degree-$d$ forms
\[
 S(s)=\sum_{\alpha\ge0}s^\alpha S_\alpha,
 \qquad S_0\ne0,
\]
satisfying the required multiplicities for general $s$.  For each
nonzero $S_\alpha$, factor out its largest power of $G_B$:
\[
 S_\alpha=G_B^{j_\alpha}R_\alpha,
 \qquad
 0\le j_\alpha\le\lfloor d/k\rfloor.
\]
Then $R_\alpha$ has degree $d-kj_\alpha$ and $R_\alpha|_B\ne0$,
since the homogeneous ideal of $B$ is generated by $G_B$.
Set
\[
 u=\min_{S_\alpha\ne0}(\alpha+j_\alpha),
 \qquad
 F(w)=\sum_{\alpha+j_\alpha=u}
       (R_\alpha|_B)w^{j_\alpha}.
\]
The sum is nonempty, and its terms have distinct exponents $j_\alpha$
because $\alpha+j_\alpha=u$.  Thus $F$ is nonzero, with algebraic
coefficients of the types in \eqref{eq:normal-polynomial}.

It remains to identify this polynomial as the normal limit and pass
the multiplicities to it.  Near a marked point choose a frame $e$ for
$\mathcal O_{\Pp^2}(1)$, put $t=G_B/e^k$, and extend a local coordinate
on $B$ to a holomorphic coordinate $\zeta$ off $B$ so that $(\zeta,t)$
are surface coordinates.  Write $r_\alpha(\zeta,t)$ for the scalar
expression of $R_\alpha$ in the frame $e^{d-kj_\alpha}$.  Substituting
$t=sw$ in the local expression of $S(s)$ and dividing by $s^u$ gives
\[
 \Phi(s,\zeta,w)
 =\sum_{S_\alpha\ne0}
 s^{\alpha+j_\alpha-u}w^{j_\alpha}
 r_\alpha(\zeta,sw).
\]
Every exponent of $s$ is nonnegative, so $\Phi$ extends
holomorphically to $s=0$ with value $F$.  Terms with
$\alpha+j_\alpha>u$ vanish there; in the remaining terms only
$r_\alpha(\zeta,0)$ contributes.

The limit respects the normal-bundle and value-bundle frames.  Indeed,
if $e'=a(\zeta,t)e$, then $t'=a(\zeta,t)^{-k}t$ and the scalar
expression of the section is multiplied by $a(\zeta,t)^{-d}$.
After $t=sw$, the transitions at $s=0$ are
\[
 w'=a(\zeta,0)^{-k}w,
 \qquad
 F'=a(\zeta,0)^{-d}F.
\]
These are the transitions of $M$ and the pullback of $L^d$.
Changing the extension of $\zeta$ off $B$ reduces at $s=0$ to its
coordinate change on $B$, since the other terms contain $t=sw$.
Thus the local rescaling realizes the intrinsic polynomial $F$ above.

Finally fix an index $i$ and use such coordinates near $a_i$.  The
functions
\[
 \zeta_i(s)=\zeta(p_i(s)),
 \qquad
 w_i(s)=\frac{t(p_i(s))}{s}
\]
extend holomorphically to $s=0$, with values representing $a_i$ and
$c_i$.  Translate the divided family to the lifted moving point:
\[
 \Psi_i(s,\xi,\omega)
   =\Phi\bigl(s,\zeta_i(s)+\xi,w_i(s)+\omega\bigr).
\]
For each general nonzero $s$, the substitution $t=sw$ is a local
biholomorphism of the surface, and multiplication by $s^{-u}$ is
multiplication by a nonzero scalar.  The multiplicity assumption on
$S(s)$ therefore implies
\[
 \partial_\xi^a\partial_\omega^b
 \Psi_i(s,0,0)=0
 \qquad(a,b\ge0,\ a+b<m).
\]
Each expression on the left is holomorphic at $s=0$ and vanishes
for general nonzero $s$, hence also at zero.  These are exactly the
order-$m$ conditions for $F$ at $(a_i,c_i)$.
\end{proof}

\subsection{The square case}

We use two basic facts about smooth projective curves.  A nonzero
section of a line bundle has an effective zero divisor whose degree
is the degree of the bundle.  In particular, a degree-zero line
bundle has a nonzero section only if it is trivial.  Also, for a
smooth complex projective curve of positive genus, $\Pic^0(B)$ is
a complex torus; its subgroup of elements of any fixed finite order
is finite, so its torsion subgroup is countable~\cite[Tag~0C1Z]{Stacks}.

\begin{proposition}[Strict inequality for square $r$]
\label{prop:squares}
Let $r=k^2$ with $k\ge4$.  There are no positive integers $d,m$ such
that $d\le km$ and $\mathsf U_r(d,m)$ holds.
\end{proposition}

\begin{proof}
Suppose that $d\le km$ and $\mathsf U_r(d,m)$ holds.  We choose
the marked points to exclude normal polynomials of fiber degree below
$m$, and the displacements to exclude fiber degree $m$.
Take a smooth degree-$k$ plane curve $B$, for example a Fermat curve,
and retain $L=\mathcal O_B(1)$ and $M=L^k$.  Its genus is
$(k-1)(k-2)/2>0$, and $\deg M=k^2=r$.  Choose distinct points
$a_1,\ldots,a_r$ so that, with $P=\sum_i a_i$, the degree-zero bundle
\[
 A=M(-P)
\]
is nontorsion.  Such a choice exists: fix $r-1$ distinct points and
vary the last one.  The map $a\mapsto\mathcal O_B(a)$ is injective
on a positive-genus curve, since a linear equivalence between two
distinct single points would give a nonconstant meromorphic function
with one simple pole, hence a degree-one map to $\Pp^1$, making the
curve rational.  Thus varying the last point gives uncountably many
distinct bundles $A$, while the torsion bundles are countable; the
finitely many prohibited coincidences may also be avoided.

The restriction sequence
\[
 0\longrightarrow\mathcal O_{\Pp^2}
 \xrightarrow{\,G_B\,}\mathcal O_{\Pp^2}(k)
 \longrightarrow\mathcal O_B(k)\longrightarrow0
\]
and $H^1(\Pp^2,\mathcal O_{\Pp^2})=0$~\cite[Tag~01XT]{Stacks} give
\[
 h^0(B,M)=\binom{k+2}{2}-1
         =\frac{k^2+3k}{2}<k^2=r.
\]
Hence one can choose
$c=(c_i)_i\in\bigoplus_i M_{a_i}$ outside the image of the evaluation
map
\[
 H^0(B,M)\longrightarrow\bigoplus_{i=1}^r M_{a_i}.
\]

Apply Proposition~\ref{prop:normal-polynomial} with these choices.  Let $j$
be the largest exponent with $f_j\ne0$.  Then
$j\le\lfloor d/k\rfloor\le m$.  If $j<m$, take $j$ derivatives with
respect to a local fiber coordinate.  Since $j$ is the largest
exponent, the result is $j!f_j$, independent of the fiber coordinate.
Differentiation lowers order by at most $j$, so $f_j$ has base
vanishing order at least $m-j$ at each $a_i$.  Its bundle has degree
\[
 \deg L^{d-kj}=k(d-kj)\le k^2(m-j)=r(m-j).
\]
Strict inequality is incompatible with these required zeros.  Equality
forces $d=km$, and removing the zeros then gives a nonzero section of
\[
 L^{k(m-j)}\bigl(-(m-j)P\bigr)=A^{m-j}.
\]
This is a nontrivial degree-zero bundle, because $A$ is nontorsion and
$m-j>0$, so it too has no nonzero section.  Thus $j<m$ is impossible.

It follows that $j=m$ and $d=km$.  In this case $f_m$ is a nonzero
constant.  At each marked fiber the restriction of $F$ is a degree-$m$
polynomial, with leading coefficient $f_m$, having a zero of order at
least $m$ at $c_i$.  In any compatible frame it therefore equals
\[
 F|_{M_{a_i}}(w)=f_m(w-c_i)^m.
\]
Comparing the coefficient of $w^{m-1}$ gives
\[
 (f_{m-1}(a_i))_i=-m f_m(c_i)_i.
\]
But $f_{m-1}\in H^0(B,M)$, and multiplication by the nonzero scalar
$-m f_m$ preserves the image of the evaluation map.  This contradicts
the choice of $c$ and finishes the proof.
\end{proof}

Together with Lemma~\ref{lem:universal}, the proposition proves
Theorem~\ref{thm:main} for square $r$, including the strict inequality
and unequal multiplicities.

\subsection{Nine zero displacements on a cubic}

For the remaining argument we will apply
Proposition~\ref{prop:normal-polynomial} with $k=3$ and with zero
normal displacements at nine distinct points of a smooth plane cubic.
Write $P=a_1+\cdots+a_9$.  Extend the coefficient list by $f_j=0$
for $j>\lfloor d/3\rfloor$.  At each $(a_i,0)$, expansion in the
fiber variable shows that multiplicity at least $m$ implies
\begin{equation}
\label{eq:nine-zeros}
 f_j\in H^0\bigl(B,L^{d-3j}(-(m-j)P)\bigr)
 \qquad(0\le j<m).
\end{equation}
Indeed, in a local base coordinate $\zeta$ centered at $a_i$, the
coefficient of $\zeta^a w^j$ must vanish whenever $a+j<m$.
Thus $f_j$ has base order at least $m-j$ at every point of $P$.
The coefficient bundles in \eqref{eq:nine-zeros} all have degree
$3(d-3m)$, independent of $j$.  These conditions hold regardless of
the additional points and normal displacements chosen in the proposition.

\begin{lemma}[Excluding degree at most $3m$]
\label{lem:cubic-low-degree}
If $r\ge10$ and $\mathsf U_r(d,m)$ holds for positive integers $d,m$,
then $d>3m$.
\end{lemma}

\begin{proof}
Suppose first that $d<3m$.  Choose a smooth plane cubic and nine
distinct points with zero normal displacements, as above, completing
them to $r$ distinct base points with arbitrary further displacements.
Every exponent occurring in \eqref{eq:normal-polynomial} is less than
$m$.  For each such exponent, \eqref{eq:nine-zeros} prescribes at
least $9(m-j)$ zeros, whereas
\[
 \deg L^{d-3j}=3(d-3j)<9(m-j).
\]
Thus every coefficient is zero, contradicting the nonzero polynomial
provided by Proposition~\ref{prop:normal-polynomial}.

Now suppose that $d=3m$.  A smooth plane cubic has genus one.  Choose
the nine points so that $A=L^3(-P)$ is nontorsion, by the same
argument used in the square case.  Complete them to $r$ distinct
base points, choosing a nonzero vector $c_{10}\in M_{a_{10}}$;
this is possible because $r\ge10$.  Apply
Proposition~\ref{prop:normal-polynomial} with these choices and zero
displacements at the first nine points.
For every $j<m$, the bundle in
\eqref{eq:nine-zeros} is $A^{m-j}$, a nontrivial degree-zero bundle.
Hence these coefficients vanish, and the nonzero normal polynomial
must be $F(w)=f_mw^m$ for a nonzero constant $f_m$.  This is nonzero at
$(a_{10},c_{10})$, so it cannot have multiplicity at least $m\ge1$
there.  This is again a contradiction.
\end{proof}
